\section{Implications if the hypotheses hold}
\label{sec:implications}

\subsection{Governance can expand while processing labor contracts}

A gate does not have to disappear for its economics to change. If verified checks become cheaper, the
enterprise can perform more of them at the same budget: more projects reviewed, more conditions
traced to closure, and more frequent reassessment. The paradox is that a stricter DGF can require
less routine case-processing labor. The gain depends on evidence quality and total assurance cost;
the relevant choice is how to use released capacity, not whether every formal review is removed.

This implication is stronger than a claim that agents write reports faster. The capacity frontier
changes the set of control programs the enterprise can afford. A firm can purchase broader coverage,
faster decisions, or more expert attention to difficult cases. These are alternative uses of the same
resource gain and cannot all be counted as independent financial savings.

\subsection{Upstream interfaces can outweigh local headcount}

The largest department need not be the most valuable first intervention. An acquisition's
inherited-system evidence or a new project's architecture-and-permission record can affect several
specialist populations simultaneously. Its value comes from the downstream work it changes, including
avoided returns. Conversely, an upstream step with few consumers or expensive reconciliation may
have little leverage.

If Hypothesis~H2 holds and measured portfolio interactions are favorable, prioritization should follow
verified dependency effects, not only local employee counts. The appropriate comparison is the full
counterfactual in Equation~\eqref{eq:pipeline}. This is upstream leverage, not a rule that earlier is
always better or that every cascade reinforces itself.

\subsection{Departments become control planes}

In the exception-routed configuration, some specialist effort moves from individual decisions to
the system that produces decisions. Procurement maintains supplier-evidence contracts; Legal
maintains approved clause interpretations and escalation rules; architects maintain reusable
patterns; Security maintains control tests and adversarial evaluations; IT production maintains recovery checks and execution limits; Compliance maintains the mapping of
rules to controls. Humans still handle novel cases and own consequential authorizations
where retained.

A \emph{control plane} is this layer of policy, evaluation, monitoring, and delegated authority over
the case-processing layer. Its work is precisely the $B_A$ and shared assurance that a credible
business case must count. The organizational shift is from applying expertise afresh to every case
toward encoding, testing, and revising the conditions under which execution is delegated.
Bainbridge's skill-retention problem remains relevant: a small team cannot safely oversee rare
failures without opportunities to maintain operational competence \citep{bainbridge1983}.

\subsection{Firm boundaries and gate redesign remain choices}

Lower coordination costs, internal or external, can move the boundary of the firm in either direction
(Section~\ref{sec:positioning}). Companies may also combine gates, remove duplicate reviews, or create
a shared evidence service rather than automate the historical sequence. The DGF is an initial
architecture to measure, not a requirement to preserve every inherited handoff. This is compatible with
Law~3: what must be preserved is every obligation and the logical order of authorizations, not each
inherited box. Inside those choices,
the gate model identifies what must be measured: whose work is removed, whose work is added, what is
verified, and where authority resides.
